\begin{table}[h!]
    \centering
    \renewcommand{\arraystretch}{1.25} % Adjust the cell height
    \setlength{\arrayrulewidth}{1pt} % Adjust the border thickness
    \begin{tabular}{p{0.3\textwidth} p{0.7\textwidth}}    
    \toprule
    \textbf{Élément} & \textbf{Description} \\
    \midrule
    ID & CU3 \\ \hline
    Cas d'utilisation & S'authentifier \\ \hline
    Description brève & L'utilisateur importe des vidéos. \\ \hline
    Acteur & Utilisateur \\ \hline
    Précondition & Authentification de l'utilisateur \\ \hline
    Scénario nominal & 
        \textbf{1.} L'utilisateur accède à la page d'importation de vidéos de site. \newline
        \textbf{2.} L'utilisateur clique sur le bouton "importer une vidéo". \newline
        \textbf{3.} L'utilisateur sélectionne la vidéo qu'il souhaite importer depuis son ordinateur. \newline
        \textbf{4.} Le système vérifie le format et la validité de la vidéo. \newline
        \textbf{5.} Le système importe la vidéo. \\ \hline
    Postcondition & Vidéo importée et affichée dans l'interface \\ \hline
    Enchaînement alternatif & \textbf{4.1} La vidéo est invalide : \newline 
        \textbullet{} \textbf{4.1.1} Le système affiche un message d'erreur indiquant le problème. \newline
        \textbullet{} \textbf{4.1.2} L'utilisateur est redirigé à l'étape 2 du scénario nominal et peut importer une nouvelle vidéo. \\
    \bottomrule
    \end{tabular} 
    \caption{Description textuelle du cas d’utilisation « Importer des vidéos »}
    \label{tab:t4}
\end{table}
